% !PRINTED 404
\weberpage{404}

Die Frage, auf die wir hier gef\"uhrt werden, ist also die, wie
man die L\"osung einer algebraischen Gleichung, deren Co\"efficienten
 von einer gewissen Anzahl von Variablen abh\"angen, auf
eine Gleichung mit einer m\"oglichst geringen Anzahl von Parametern
, und insbesondere, unter welchen Umst\"anden und mit
welchem H\"ulfsmittel man sie auf Gleichungen mit nur einem
Parameter zur\"uckf\"uhren kann.

Wir stellen uns demnach jetzt die Frage, unter welchen
Voraussetzungen bei einer \gesperrt{allgemeinen Gleichung $n^{\text{ten}}$ Grades
Resolventen mit nur einem Parameter existiren.}

Es sei $x_0, x_1, \ldots, x_{n-1}$ ein System von $n$ unabh\"angigen Variablen,
\weq{1}{u = \Phi(x_0, x_1, \ldots, x_{n-1})}
eine rationale Function dieser Variablen, die durch die Permutationen
 der \gesperrt{alternirenden Gruppe} der $n$ Buchstaben $x$ in die
von einander verschiedenen Functionen
\weq{2}{u,\ u_1, u_2, \ldots u_{\nu-1}}
\"ubergeht, worin $\nu$ gleich oder kleiner als der Grad der alternirenden
 Gruppe sein kann, und $z$ eine Function der Variablen $x$
ist, die durch die alternirende Gruppe unge\"andert bleibt. Es
fragt sich, wann eine rationale Gleichung $\nu^{\text{ten}}$ Grades in Bezug
auf $u$ besteht,
\weq{3}{F(u, z) = 0,}
die durch die $\nu$ Functionen (2) identisch befriedigt wird? Die
Co\"efficienten in (3) m\"ussen von den $x$ unabh\"angig und also reine
Zahlen sein. Eine Beschr\"ankung des Rationalit\"atsbereiches im
Gebiete der Zahlen lassen wir einstweilen nicht eintreten.

Die Gleichung (3) ist, wenn $\nu > 1$ ist, eine Resolvente der
Gleichung $n^{\text{ten}}$ Grades, deren Wurzeln die $x$ sind, die nur von dem
\gesperrt{einen Parameter} $z$ abh\"angt. Den Fall $\nu = 1$ schliessen wir aus.

\begin{center}
\S.\,107.

\gesperrt{Satz von L\"uroth.}
\end{center}

Es ist zun\"achst folgender H\"ulfssatz zu beweisen \textsuperscript{1}):

1. \gesperrt{Ist $u, u_1, u_2, \ldots, u_{\nu-1}$ ein System rationaler Functionen
 einer Variablen $t$, so giebt es eine rationale
Function der $u, u_1, \ldots, u_{\nu-1}$:}
\[
\vartheta = \chi(u, u_1, \ldots, u_{\nu-1})
\]

\begin{center}
\rule{0.3\linewidth}{0.4pt}
\end{center}

\textsuperscript{1}) \gesperrt{L\"uroth}, Mathematische Annalen, Bd. IX.
%CONT>>
